%==============================================================================
%  BAB I - LANDASAN TEORI
%==============================================================================
\section{Landasan Teori}

Milestone 4 merupakan tahap akhir dari pengembangan kompiler bahasa Arion. Pada
tahap ini, \emph{Decorated AST} yang dihasilkan Milestone 3 diubah menjadi
\emph{Intermediate Code} berbentuk \emph{Three-Address Code} (TAC), lalu
dieksekusi oleh sebuah \emph{interpreter} berbasis \emph{stack machine} hingga
menghasilkan keluaran nyata program. Bab ini menjelaskan teori dan konsep yang
melandasi kedua komponen tersebut.

%------------------------------------------------------------------------------
\subsection{Arsitektur Kompiler: Front-end dan Back-end}

Kompiler modern umumnya dibagi menjadi dua fase arsitektur: \textbf{front-end}
dan \textbf{back-end}. Milestone 1 sampai 3 (analisis leksikal, sintaksis, dan
semantik) merupakan bagian front-end yang berurusan langsung dengan teks kode
sumber---memotongnya menjadi token, memvalidasi tata bahasanya, lalu memastikan
maknanya benar. Keluaran front-end adalah struktur data murni berupa
\emph{Decorated AST} yang sudah bebas dari teks mentah.

Milestone 4 sepenuhnya berada pada area back-end. Fase ini tidak lagi peduli
pada detail sintaksis Arion (titik koma, kurung, atau spasi), melainkan berfokus
pada dua tugas inti: \textbf{sintesis}, yaitu mengonversi Decorated AST menjadi
instruksi operasional standar (Intermediate Code), dan \textbf{eksekusi}, yaitu
mengelola memori serta menjalankan instruksi tersebut pada mesin virtual. Pemisahan
yang tegas ini memberi keuntungan rekayasa: bila suatu saat sintaks Arion diubah,
cukup front-end yang dirombak selama Decorated AST yang dihasilkan tetap sama,
back-end tidak perlu disentuh.

\begin{figure}[H]
\centering
\begin{tikzpicture}[
  node distance=0.55cm and 0.9cm,
  kotak/.style={draw, rounded corners, minimum height=1cm, minimum width=2.3cm,
    align=center, font=\footnotesize, fill=biruITB!8},
  kotakBE/.style={draw, rounded corners, minimum height=1cm, minimum width=2.3cm,
    align=center, font=\footnotesize, fill=hijauJudul!12},
  panah/.style={-{Stealth[length=2.5mm]}, thick}
]
  \node[kotak] (src) {Source\\Code};
  \node[kotak, right=of src] (tok) {Tokens};
  \node[kotak, right=of tok] (pt) {Parse\\Tree};
  \node[kotak, right=of pt] (ast) {Decorated\\AST};
  \node[kotakBE, below=1.1cm of ast] (tac) {Intermediate\\Code (TAC)};
  \node[kotakBE, left=of tac] (vm) {Stack\\Machine};
  \node[kotakBE, left=of vm] (out) {Program\\Output};

  \draw[panah] (src) -- (tok);
  \draw[panah] (tok) -- (pt);
  \draw[panah] (pt) -- (ast);
  \draw[panah] (ast) -- (tac);
  \draw[panah] (tac) -- (vm);
  \draw[panah] (vm) -- (out);

  \node[font=\scriptsize\itshape, above=0.05cm of tok] {Front-end (M1--M3)};
  \node[font=\scriptsize\itshape, below=0.05cm of vm] {Back-end (M4)};
\end{tikzpicture}
\caption{Pipeline kompiler Arion. Milestone 4 menambahkan tahap pembentukan
Intermediate Code dan eksekusi pada stack machine.}
\label{fig:pipeline}
\end{figure}

%------------------------------------------------------------------------------
\subsection{Intermediate Code Generation}

\subsubsection{Decorated AST sebagai Sumber Kebenaran}

\emph{Decorated AST} adalah \emph{Abstract Syntax Tree} yang setiap simpulnya
telah dianotasi (didekorasi) dengan informasi semantik, terutama \textbf{tipe
data} hasil ekspresi dan \textbf{referensi tabel simbol} untuk variabel. Anotasi
ini bersifat wajib bagi generator kode. Sebagai ilustrasi, ketika generator
menemui ekspresi \texttt{a + b}, tanpa informasi tipe ia tidak dapat memutuskan
apakah operasi tersebut penjumlahan bilangan atau penggabungan string. Dengan
Decorated AST, generator dapat memilih instruksi Intermediate Code yang tepat.

\subsubsection{Three-Address Code (TAC)}

Langkah berikutnya adalah meratakan (\emph{flattening}) Decorated AST menjadi
daftar instruksi linier. Format yang digunakan adalah \textbf{Three-Address
Code}, dinamai demikian karena tiap instruksi maksimal melibatkan tiga alamat
(dua operand sumber dan satu tujuan), dengan struktur dasar
\texttt{hasil = operand1 operator operand2}.

Karena CPU tidak dapat menghitung ekspresi panjang sekaligus, TAC memecah ekspresi
kompleks menggunakan \textbf{variabel sementara} (\emph{temporary}, ditulis
\texttt{t1}, \texttt{t2}, \dots). Pohon ekspresi dievaluasi dari cabang terdalam
menuju akar. Misalnya, ekspresi \texttt{x := (a + b) * (c - d)} diratakan menjadi:

\begin{lstlisting}[style=kodeTAC]
t1 := a + b      { evaluasi tanda kurung pertama }
t2 := c - d      { evaluasi tanda kurung kedua   }
t3 := t1 * t2    { kalikan hasil keduanya        }
x  := t3         { simpan hasil akhir ke x       }
\end{lstlisting}

Pada implementasi berbasis stack machine, peran variabel sementara digantikan oleh
puncak stack: hasil antara cukup ditinggalkan di atas stack untuk dikonsumsi
operasi berikutnya, sehingga tidak perlu dialokasikan nama \texttt{t1}, \texttt{t2}
secara eksplisit.

\subsubsection{Translasi Control Flow}

Tantangan terbesar pada pembentukan Intermediate Code adalah meratakan struktur
\emph{control flow} bersarang (\texttt{if-else}, \texttt{while}) menjadi daftar
linier. Solusinya menggunakan kombinasi \textbf{Label} dan \textbf{Jump}:

\textbf{Label} adalah penanda statis posisi suatu baris dalam TAC.
\textbf{Jump mutlak} (\texttt{JMP}) melompat tanpa syarat ke sebuah label,
sedangkan \textbf{jump bersyarat} (\texttt{JPC}) melompat hanya jika kondisi yang
dievaluasi bernilai \emph{false}. Konvensi melompat-jika-salah ini disebut
\emph{IF\_FALSE} dan memanfaatkan prinsip \emph{fall-through}: ketika kondisi
benar, eksekusi langsung berlanjut ke instruksi berikutnya tanpa lompatan,
sehingga alur kontrol menjadi lebih sederhana. Pada tugas besar ini konvensi
\emph{IF\_FALSE} bersifat wajib.

Berikut pola translasi \texttt{if-else} (kiri) dan \texttt{while} (kanan) menjadi
Intermediate Code.

\begin{center}
\begin{minipage}[t]{0.46\textwidth}
\centering\textbf{IF-ELSE}
\begin{lstlisting}[style=kodeTAC, numbers=none]
  t1 := x > 10
  JPC t1 -> L_ELSE
  y := 1
  JMP L_END
L_ELSE:
  y := 0
L_END:
\end{lstlisting}
\end{minipage}\hfill
\begin{minipage}[t]{0.46\textwidth}
\centering\textbf{WHILE}
\begin{lstlisting}[style=kodeTAC, numbers=none]
L_START:
  t1 := i < 5
  JPC t1 -> L_END
  i := i + 1
  JMP L_START
L_END:
\end{lstlisting}
\end{minipage}
\end{center}

%------------------------------------------------------------------------------
\subsection{Katalog Instruksi Intermediate Code}

Intermediate Code Arion menggunakan himpunan instruksi bergaya ORKOM. Tiap
instruksi terdiri atas \emph{mnemonic} dan operand opsional. Tabel~\ref{tab:instr}
merangkum instruksi yang didukung, sedangkan Tabel~\ref{tab:opr} merinci operasi
yang dipanggil melalui instruksi \texttt{OPR}.

\begin{table}[H]
\centering
\caption{Daftar instruksi Intermediate Code.}
\label{tab:instr}
\renewcommand{\arraystretch}{1.15}
\begin{tabularx}{\textwidth}{@{}l l X@{}}
\toprule
\textbf{Instruksi} & \textbf{Operasi} & \textbf{Keterangan} \\
\midrule
\texttt{LIT v} & Load Literal     & Memasukkan nilai literal \texttt{v} ke stack. \\
\texttt{LOD a} & Load Value       & Memuat nilai dari alamat \texttt{a} ke stack. \\
\texttt{STO a} & Store Value      & Mengambil nilai puncak stack, menyimpannya ke alamat \texttt{a}. \\
\texttt{CAL l} & Call             & Memanggil fungsi/prosedur pada baris \texttt{l}. \\
\texttt{INT m} & Initiate Memory  & Mengalokasikan memori (stack frame) berukuran \texttt{m}. \\
\texttt{JMP l} & Unconditional Jump & Lompat ke baris \texttt{l} tanpa syarat. \\
\texttt{JPC l} & Conditional Jump & Lompat ke baris \texttt{l} bila kondisi puncak stack salah. \\
\texttt{OPR o} & Operation        & Menjalankan operasi bernomor \texttt{o} (lihat Tabel~\ref{tab:opr}). \\
\texttt{RET}   & Return           & Keluar dari fungsi/prosedur atau mengakhiri program. \\
\bottomrule
\end{tabularx}
\end{table}

\begin{table}[H]
\centering
\caption{Daftar operasi instruksi \texttt{OPR}.}
\label{tab:opr}
\renewcommand{\arraystretch}{1.1}
\small
\begin{tabularx}{\textwidth}{@{}c l X | c l X@{}}
\toprule
\textbf{No} & \textbf{Nama} & \textbf{Fungsi} & \textbf{No} & \textbf{Nama} & \textbf{Fungsi} \\
\midrule
1 & NEG & Negasi puncak stack          & 8  & NEQ   & Tidak sama dengan \\
2 & ADD & Penjumlahan                  & 9  & LSS   & Kurang dari \\
3 & SUB & Pengurangan                  & 10 & GEQ   & Lebih dari sama dengan \\
4 & MUL & Perkalian                    & 11 & GTR   & Lebih dari \\
5 & DIV & Pembagian                    & 12 & LEQ   & Kurang dari sama dengan \\
6 & MOD & Modulus                      & 13 & WRT   & Cetak output \\
7 & EQL & Sama dengan                  & 14 & WRTLN & Cetak output + newline \\
\bottomrule
\end{tabularx}
\end{table}

%------------------------------------------------------------------------------
\subsection{Arsitektur Mesin Eksekusi: Interpreter dan Stack}

\subsubsection{Virtual Machine dan Siklus Eksekusi}

\emph{Interpreter} pada Milestone 4 berperan sebagai sebuah \emph{Virtual Machine}
(VM). Sama seperti CPU sungguhan, VM membaca instruksi satu per satu menggunakan
penunjuk \emph{Instruction Pointer} (IP) atau \emph{Program Counter}. VM bekerja
dalam \textbf{siklus eksekusi} tiga tahap: \textbf{Fetch} (membaca instruksi yang
ditunjuk IP), \textbf{Decode} (membedah mnemonic dan operand-nya), dan
\textbf{Execute} (melakukan aksi nyata). Setelah Execute, IP bergerak maju ke
instruksi berikutnya, kecuali pada instruksi lompat (\texttt{JMP}/\texttt{JPC})
yang menggeser IP langsung ke baris tujuan. Siklus berulang hingga IP mencapai
akhir program atau menemui \texttt{RET}.

\subsubsection{Memori Eksekusi: The Stack}

Struktur memori utama VM adalah \textbf{stack} yang bekerja dengan prinsip
\textbf{LIFO} (\emph{Last In, First Out})---data hanya bisa dimasukkan
(\emph{push}) dan diambil (\emph{pop}) dari puncak. Stack dipilih, bukan array
biasa, karena berkaitan erat dengan \emph{scope} dan pemanggilan fungsi. Setiap
kali program memasuki blok atau memanggil fungsi, VM membuat \emph{stack frame}
baru. Pada konvensi yang dipakai, tiga sel pertama setiap frame selalu terisi
\textbf{Static Link} (alamat 0), \textbf{Dynamic Link} (alamat 1), dan
\textbf{Return Address} (alamat 2); sel selanjutnya (mulai alamat 3) digunakan
untuk variabel. Akses variabel dilakukan relatif terhadap \emph{Base Pointer}
(BP) frame aktif, sementara \emph{Stack Pointer} (SP) menunjuk puncak stack.

\begin{figure}[H]
\centering
\begin{tikzpicture}[
  sel/.style={draw, minimum width=5.5cm, minimum height=0.7cm, font=\footnotesize, anchor=south west},
  selV/.style={sel, fill=hijauJudul!12},
  selR/.style={sel, fill=biruITB!10},
]
  \node[selR] at (0,0)   (a0) {alamat 0 -- Static Link};
  \node[selR] at (0,0.7) (a1) {alamat 1 -- Dynamic Link};
  \node[selR] at (0,1.4) (a2) {alamat 2 -- Return Address};
  \node[selV] at (0,2.1) (a3) {alamat 3 -- Variable (x)};
  \node[selV] at (0,2.8) (a4) {alamat 4 -- Variable (y)};
  \node[sel, fill=white] at (0,3.5) (a5) {puncak -- operand/hasil sementara};

  \node[font=\footnotesize, right=0.3cm of a0] {$\leftarrow$ BP (basis frame)};
  \node[font=\footnotesize, right=0.3cm of a5] {$\leftarrow$ SP (puncak stack)};
\end{tikzpicture}
\caption{Struktur stack frame: tiga sel administratif diikuti sel variabel; ruang
di atasnya dipakai sebagai operand dan hasil sementara perhitungan.}
\label{fig:frame}
\end{figure}

%------------------------------------------------------------------------------
\subsection{Keamanan Runtime dan Penanganan Kerentanan}

Interpreter yang tangguh tidak hanya menjalankan kode yang benar, tetapi juga
mampu bertahan (tidak \emph{crash}) ketika diberi kode yang salah atau tidak
masuk akal. Berbeda dengan bahasa tingkat rendah yang menyerahkan keamanan kepada
\emph{programmer}, bahasa yang berjalan di atas VM wajib memiliki mekanisme
perlindungan internal. Sebagai bagian dari spesifikasi (termasuk tantangan bonus),
interpreter Arion mendeteksi dan menangani sejumlah kerentanan berikut.

\textbf{Kerentanan pada stack.} \emph{Stack overflow} terjadi ketika frame terus
ditambahkan melampaui batas memori, umumnya akibat rekursi tanpa \emph{base case};
interpreter membatasi kedalaman frame dan melempar error bila terlampaui.
\emph{Stack underflow} terjadi ketika operasi \emph{pop} dipanggil saat stack
kosong. \emph{Stack smashing} adalah perusakan \emph{return address} oleh data yang
meluber, dideteksi menggunakan \emph{canary} pada sel return address.
\emph{Stack corruption} terjadi ketika jumlah \emph{push} dan \emph{pop} tidak
simetris sehingga indeks frame bergeser; ini dideteksi dengan memverifikasi posisi
stack pointer saat keluar dari frame.

\textbf{Kerentanan memori dan alur.} \emph{Out-of-Bounds Access} dicegah dengan
memvalidasi setiap alamat pada \texttt{LOD}/\texttt{STO} terhadap batas memori.
\emph{Invalid Jump Target} dicegah dengan memverifikasi target \texttt{JMP},
\texttt{JPC}, dan \texttt{CAL} berada dalam rentang program sebelum IP dipindahkan.

\textbf{Kerentanan aritmatika.} Tipe \emph{integer} memiliki kapasitas terbatas;
operasi yang melampaui batas atas (\emph{overflow}) atau bawah (\emph{underflow})
akan menyebabkan \emph{wrap-around} yang mengacaukan perhitungan. Interpreter
mendeteksi kondisi ini sebelum hasil sempat membungkus, lalu melempar
\textbf{OverflowError} atau \textbf{UnderflowError}. Demikian pula pembagian dengan
nol dicegat sebagai \textbf{DivisionByZeroError}.

\textbf{Runtime type checking.} Karena Arion bersifat \emph{strongly-typed} dan
sebagian besar tipe telah diverifikasi pada analisis semantik, runtime type
checking tidak diwajibkan. Meski demikian, interpreter tetap memvalidasi bahwa
operand operasi aritmatika benar-benar bernilai numerik saat eksekusi; bila tidak,
ia melempar \textbf{TypeMismatchError} alih-alih melakukan konversi diam-diam.
